\providecommand{\GL}{\operatorname{GL}}
\providecommand{\SO}{\operatorname{SO}}
\providecommand{\Orth}{\operatorname{O}}
\providecommand{\Sym}{\operatorname{Sym}}
\providecommand{\Stab}{\operatorname{Stab}}
\providecommand{\Orb}{\operatorname{Orb}}
\providecommand{\Lie}[1]{\mathfrak{#1}}

\chapter{Fundamentos Matem\'aticos}
\label{ch:foundations}

La premisa del aprendizaje profundo geom\'etrico, planteada en el
\Cref{ch:introduction}, es que la estructura geom\'etrica de los datos debe
dictar la estructura del modelo. Para volver operativa esta premisa necesitamos
un lenguaje preciso para dos clases de estructura: las \emph{simetr\'ias} que
posee un problema y los \emph{espacios} en los que viven los datos. Este
cap\'itulo desarrolla ese lenguaje. Comenzamos con los grupos y su forma de
actuar sobre conjuntos, que formalizan la simetr\'ia; pasamos despu\'es a las
representaciones lineales, el mecanismo por el cual una simetr\'ia abstracta
act\'ua sobre los espacios de caracter\'isticas de una red neuronal; describimos
los espacios homog\'eneos y los cocientes, que unifican dominios tan distintos
como la esfera y el plano bajo una misma construcci\'on; y cerramos con las
variedades y los espacios tangentes, la geometr\'ia sobre la que se asientan los
cap\'itulos posteriores. El tratamiento es deliberadamente autocontenido pero
selectivo: enunciamos las definiciones y teoremas que se emplean en el resto del
libro y demostramos solo aquellos resultados cuyas pruebas son ilustrativas.
Referencias est\'andar para este material son~\citet{lee2012introduction} para
variedades, \citet{munkres2000topology} para topolog\'ia y
\citet{Kirillov2008,Fulton1991} para grupos de Lie y teor\'ia de
representaciones.

\section{Grupos y acciones de grupo}
\label{sec:found:groups}

Una simetr\'ia es una transformaci\'on que deja invariable una magnitud de
inter\'es. El conjunto de todas esas transformaciones nunca es arbitrario: la
composici\'on de dos simetr\'ias es de nuevo una simetr\'ia, toda simetr\'ia se
puede deshacer y no hacer nada es una simetr\'ia. Estos tres hechos son
exactamente los axiomas de un grupo.

\begin{definition}[Grupo]
\label{def:found:group}
Un \emph{grupo} es un conjunto $\group$ dotado de una operaci\'on binaria
$\cdot : \group \times \group \to \group$ que satisface:
\begin{enumerate}
  \item \textbf{Asociatividad:} $(a \cdot b) \cdot c = a \cdot (b \cdot c)$ para
    todo $a, b, c \in \group$.
  \item \textbf{Identidad:} existe $e \in \group$ con
    $e \cdot a = a \cdot e = a$ para todo $a \in \group$.
  \item \textbf{Inversos:} para cada $a \in \group$ existe $a^{-1} \in \group$
    con $a \cdot a^{-1} = a^{-1} \cdot a = e$.
\end{enumerate}
El grupo es \emph{abeliano} (o conmutativo) si $a \cdot b = b \cdot a$ para todo
$a, b$.
\end{definition}

Las traslaciones de una imagen forman un grupo abeliano: desplazar por $(a,b)$ y
luego por $(c,d)$ da el mismo resultado en cualquier orden. Las rotaciones en
tres dimensiones, en cambio, no conmutan, y su grupo es no abeliano. Un
subconjunto $H \subseteq \group$ que es a su vez un grupo con la operaci\'on
heredada es un \emph{subgrupo}, escrito $H \le \group$. Un subgrupo $H$ es
\emph{normal} ($H \triangleleft \group$) si $gHg^{-1} = H$ para todo
$g \in \group$; la normalidad es precisamente la condici\'on para que las clases
laterales de $H$ hereden una estructura de grupo, como retomamos en la
\Cref{sec:found:homogeneous}.

El puente entre dos grupos es una aplicaci\'on que preserva sus operaciones.

\begin{definition}[Homomorfismo]
\label{def:found:homomorphism}
Una aplicaci\'on $\phi : \group \to \mathcal{H}$ entre grupos es un
\emph{homomorfismo} si $\phi(g_1 g_2) = \phi(g_1)\phi(g_2)$ para todo
$g_1, g_2 \in \group$. Un homomorfismo biyectivo es un \emph{isomorfismo},
escrito $\group \cong \mathcal{H}$. El \emph{n\'ucleo}
$\ker(\phi) = \{g : \phi(g) = e_{\mathcal{H}}\}$ es siempre un subgrupo normal, y
la \emph{imagen} $\operatorname{Im}(\phi) = \{\phi(g) : g \in \group\}$ es un
subgrupo de $\mathcal{H}$.
\end{definition}

Dos grupos isomorfos son algebraicamente indistinguibles: solo difieren en los
nombres de sus elementos. Los homomorfismos son el recurso formal mediante el
cual las simetr\'ias abstractas se vuelven transformaciones computables, y en los
\Cref{ch:priors,ch:groups} son lo que hace implementable la equivarianza.

Los grupos son abstractos; los datos son concretos. Para aplicar un grupo a
im\'agenes, grafos o nubes de puntos debemos decir c\'omo sus elementos
transforman los puntos de un conjunto. Esta es la noci\'on de acci\'on de grupo.

\begin{definition}[Acci\'on de grupo]
\label{def:found:action}
Una \emph{acci\'on} de un grupo $\group$ sobre un conjunto $X$ es una aplicaci\'on
$\group \times X \to X$, escrita $(g, x) \mapsto g \cdot x$, que satisface
\[
  (g_1 g_2) \cdot x = g_1 \cdot (g_2 \cdot x), \qquad e \cdot x = x
\]
para todo $g_1, g_2 \in \group$ y $x \in X$.
\end{definition}

Toda acci\'on equivale a un homomorfismo en el grupo sim\'etrico: la aplicaci\'on
$\tilde\rho(g)\colon x \mapsto g\cdot x$ es una biyecci\'on de $X$ (su inversa es
$\tilde\rho(g^{-1})$), y $\tilde\rho : \group \to \operatorname{Sym}(X)$ es un
homomorfismo. Estudiar c\'omo act\'ua $\group$ sobre $X$ es, por tanto, lo mismo
que estudiar homomorfismos $\group \to \operatorname{Sym}(X)$. Las acciones
relevantes para este libro son familiares:

\begin{example}[Acciones en aprendizaje profundo]
\label{ex:found:actions}
\leavevmode
\begin{itemize}
  \item Las traslaciones $\Z^2$ act\'uan sobre im\'agenes discretas desplazando
    p\'ixeles.
  \item Las rotaciones $\SO(3)$ act\'uan sobre nubes de puntos por
    multiplicaci\'on matricial.
  \item Las permutaciones $S_n$ act\'uan sobre grafos de $n$ nodos
    reetiquet\'andolos.
\end{itemize}
Son, respectivamente, los grupos de simetr\'ia de las rejillas, los datos
geom\'etricos y los grafos, y organizan los
\Cref{ch:grids,ch:groups,ch:graphs}.
\end{example}

Muchos grupos de simetr\'ia de inter\'es se construyen a partir de otros m\'as
simples. El grupo eucl\'ideo de los movimientos r\'igidos, por ejemplo, combina
traslaciones con rotaciones y reflexiones, pero no como un producto simple: una
rotaci\'on reorienta la traslaci\'on que la sigue. El \emph{producto
semidirecto} $N \rtimes_\phi H$ captura esto. Dado un homomorfismo
$\phi : H \to \operatorname{Aut}(N)$ que describe c\'omo act\'ua $H$ sobre $N$,
es el conjunto $N \times H$ con el producto
\[
  (n_1, h_1)\cdot(n_2, h_2) = \bigl(n_1\,\phi(h_1)(n_2),\; h_1 h_2\bigr).
\]
El grupo eucl\'ideo es $E(n) = \R^n \rtimes \Orth(n)$, con
$(t_1, R_1)\cdot(t_2, R_2) = (t_1 + R_1 t_2,\, R_1 R_2)$: primero se rota, luego
se traslada. Su subgrupo que preserva la orientaci\'on
$SE(n) = \R^n \rtimes \SO(n)$ es el grupo de simetr\'ia de los datos r\'igidos
tridimensionales.

\paragraph{Grupos de Lie.}
Las simetr\'ias anteriores se dividen en dos clases. Las permutaciones forman un
grupo \emph{discreto}; las rotaciones y traslaciones forman familias
\emph{continuas}. Para hacer c\'alculo en el caso continuo ---para diferenciar y
optimizar sobre transformaciones--- exigimos que el grupo sea tambi\'en una
variedad diferenciable (noci\'on que se precisa en la \Cref{sec:found:manifolds}).

\begin{definition}[Grupo de Lie]
\label{def:found:liegroup}
Un \emph{grupo de Lie} es un grupo $G$ que es a la vez una variedad diferenciable
y para el cual la multiplicaci\'on $(g,h) \mapsto gh$ y la inversi\'on
$g \mapsto g^{-1}$ son aplicaciones suaves.
\end{definition}

Los grupos de Lie m\'as importantes en este libro son los grupos matriciales. El
\emph{grupo ortogonal especial}
\[
  \SO(n) = \{ R \in \R^{n\times n} : R^\top R = I_n,\ \det R = 1 \}
\]
de las rotaciones es un grupo de Lie compacto de dimensi\'on
$\tfrac{n(n-1)}{2}$; al eliminar la condici\'on del determinante se obtiene el
grupo ortogonal $\Orth(n)$, que adem\'as contiene reflexiones. Para $n=2$ una
rotaci\'on es un \'unico \'angulo,
$R_\theta = \left(\begin{smallmatrix}\cos\theta & -\sin\theta\\ \sin\theta &
\cos\theta\end{smallmatrix}\right)$, de modo que $\SO(2)$ es el c\'irculo;
$\SO(3)$ tiene dimensi\'on tres, los tres ejes de rotaci\'on.

Un grupo de Lie es curvo, pero es lineal a primer orden. Su espacio tangente en
la identidad, $\Lie{g} = T_e G$, lleva un corchete bilineal antisim\'etrico
$[\cdot,\cdot]$ (el \emph{corchete de Lie}) que lo convierte en un
\emph{\'algebra de Lie}: un espacio vectorial con un corchete bilineal,
antisim\'etrico y que satisface la identidad de Jacobi
$[X,[Y,Z]] + [Y,[Z,X]] + [Z,[X,Y]] = 0$. El \'algebra de Lie lineariza el grupo,
y el paso de vuelta al grupo es el \emph{mapa exponencial}
$\exp : \Lie{g} \to G$, que para grupos matriciales es la exponencial de
matrices $\exp(A) = \sum_{k\ge 0} A^k/k!$. Para $\SO(2)$, con
$J = \left(\begin{smallmatrix}0 & -1\\ 1 & 0\end{smallmatrix}\right)$ generando
$\Lie{so}(2)$,
\[
  \exp(\theta J) = \begin{pmatrix}\cos\theta & -\sin\theta\\
  \sin\theta & \cos\theta\end{pmatrix} = R_\theta,
\]
de modo que la exponencial convierte un \'angulo (un elemento del \'algebra) en
una rotaci\'on (un elemento del grupo). Esta es la raz\'on por la que las
\'algebras de Lie importan en el dise\~no de redes neuronales: la optimizaci\'on
puede realizarse en el espacio vectorial plano $\Lie{g}$ con gradientes
ordinarios, y la transformaci\'on resultante aplicarse en el grupo mediante
$\exp$.

\section{Representaciones}
\label{sec:found:representations}

Una acci\'on de grupo mueve los puntos de un dominio. Dentro de una red
neuronal, sin embargo, las simetr\'ias deben actuar sobre \emph{espacios de
caracter\'isticas} ---espacios vectoriales de activaciones---. La noci\'on que
hace que una simetr\'ia act\'ue linealmente es la representaci\'on.

\begin{definition}[Representaci\'on]
\label{def:found:representation}
Una \emph{representaci\'on} de un grupo $\group$ sobre un espacio vectorial $V$
sobre un cuerpo $\mathbb{K}$ es un homomorfismo $\rho : \group \to \GL(V)$, donde
$\GL(V)$ es el grupo de aplicaciones lineales invertibles de $V$. De forma
equivalente, a cada $g$ se le asigna una matriz invertible $\rho(g)$ de modo que
$\rho(g_1 g_2) = \rho(g_1)\rho(g_2)$. Para un grupo de Lie se exige adem\'as que
$\rho$ sea suave.
\end{definition}

Las representaciones traducen simetr\'ias abstractas en operaciones lineales
concretas sobre los canales de una red. Una representaci\'on suave de un grupo de
Lie se diferencia a una representaci\'on de su \'algebra de Lie,
$d\rho : \Lie{g} \to \Lie{gl}(V)$ con
$d\rho(X) = \frac{d}{dt}\big|_{0}\,\rho(\exp(tX))$, que preserva el corchete:
$d\rho([X,Y]) = [d\rho(X), d\rho(Y)]$. Esto es de nuevo lo que permite optimizar
en el \'algebra lineal y actuar en el grupo.

La pregunta b\'asica de la teor\'ia de representaciones es c\'omo se descompone
una representaci\'on.

\begin{definition}[Representaci\'on irreducible]
\label{def:found:irreducible}
Una representaci\'on $\rho : \group \to \GL(V)$ con $V \ne \{0\}$ es
\emph{irreducible} si los \'unicos subespacios $W \subseteq V$ con
$\rho(g)W \subseteq W$ para todo $g$ son $\{0\}$ y $V$.
\end{definition}

Las representaciones irreducibles son los bloques fundamentales con los que se
construyen todas las dem\'as. Para $\SO(2)$ las irreducibles (complejas) son
unidimensionales, indexadas por $n \in \Z$:
\[
  \rho_n : \SO(2) \to \C^*, \qquad \rho_n(R_\theta) = e^{i n\theta}.
\]
Cada $n$ es una \emph{frecuencia}, y estas son exactamente las funciones base de
la serie de Fourier cl\'asica. El siguiente teorema afirma que para grupos
compactos toda representaci\'on se descompone en irreducibles.

\begin{theorem}[Reducibilidad completa]
\label{thm:found:decomposition}
Toda representaci\'on de dimensi\'on finita de un grupo compacto $\group$ se
descompone como suma directa de representaciones irreducibles,
$V = \bigoplus_i n_i V_i$, donde cada $V_i$ es irreducible y $n_i \in \mathbb{N}$
es su multiplicidad.
\end{theorem}

La demostraci\'on descansa en un \'unico recurso: promediar sobre el grupo. Los
grupos compactos poseen una \'unica medida de probabilidad invariante, la
\emph{medida de Haar} $\mu$, que satisface $\mu(gA) = \mu(Ag) = \mu(A)$ y
$\mu(\group) = 1$. Dado cualquier producto interno
$\langle\cdot,\cdot\rangle_0$ en $V$, el promedio
$\langle u, v\rangle = \int_\group \langle \rho(g)u, \rho(g)v\rangle_0\,d\mu(g)$
es un producto interno $\group$-invariante. Respecto de \'el, el complemento
ortogonal de cualquier subespacio invariante es de nuevo invariante, de modo que
siempre se puede separar un subespacio invariante; la inducci\'on sobre
$\dim V$ concluye el argumento. La medida de Haar es tambi\'en exactamente el
promediado que define la convoluci\'on de grupo, y su invarianza es lo que
garantiza la equivarianza del operador resultante ---el hecho central tras las
redes equivariantes a grupos del \Cref{ch:groups}~\citep{cohen2016group}.

Las aplicaciones entre irreducibles est\'an fuertemente restringidas, hecho
conocido como lema de Schur.

\begin{lemma}[Schur]
\label{lem:found:schur}
Sean $(\rho, V)$ y $(\sigma, W)$ representaciones irreducibles y sea
$T : V \to W$ tal que $T\rho(g) = \sigma(g)T$ para todo $g$ (un entrelazador).
Entonces $T$ es cero o un isomorfismo. Si adem\'as $V = W$ y el cuerpo es $\C$,
entonces $T = \lambda\,\mathrm{Id}$ para alg\'un escalar $\lambda$.
\end{lemma}

\begin{proof}
El n\'ucleo de $T$ es invariante ($v \in \ker T \Rightarrow T\rho(g)v =
\sigma(g)Tv = 0$), luego por la irreducibilidad de $\rho$ es $\{0\}$ o $V$;
an\'alogamente la imagen es invariante, luego por la irreducibilidad de $\sigma$
es $\{0\}$ o $W$. Por tanto $T = 0$ o $T$ es una biyecci\'on. Para la segunda
parte, sobre $\C$ la aplicaci\'on $T$ tiene un autovalor $\lambda$; entonces
$T - \lambda\,\mathrm{Id}$ es un entrelazador con n\'ucleo no trivial, de donde
se anula.
\end{proof}

El lema de Schur es la ra\'iz algebraica de la eficiencia param\'etrica de las
redes equivariantes: una capa lineal equivariante entre representaciones se ve
obligada a ser diagonal por bloques en la descomposici\'on irreducible, con cada
bloque proporcional a la identidad, de modo que el n\'umero de par\'ametros
libres lo dicta la simetr\'ia y no el tama\~no de la entrada. Esta
descomposici\'on culmina, para grupos de Lie compactos, en el teorema de
Peter--Weyl.

\begin{theorem}[Peter--Weyl]
\label{thm:found:peterweyl}
Sea $\group$ un grupo de Lie compacto. El espacio $L^2(\group)$ de funciones de
cuadrado integrable se descompone como
\[
  L^2(\group) = \bigoplus_{\lambda \in \widehat{\group}} d_\lambda\, V_\lambda,
\]
donde $\widehat{\group}$ es el conjunto de clases de equivalencia de
representaciones irreducibles, $V_\lambda$ es el espacio de la irreducible
$\lambda$, y $d_\lambda = \dim V_\lambda$.
\end{theorem}

El teorema de Peter--Weyl es la generalizaci\'on de la serie de Fourier a grupos
no abelianos: las entradas matriciales de las irreducibles desempe\~nan el papel
de las exponenciales $e^{in\theta}$ y forman una base ortonormal completa de
$L^2(\group)$. Para $\group = \SO(2)$ se reduce exactamente al desarrollo de
Fourier cl\'asico $f(\theta) = \sum_n \hat f_n\,e^{in\theta}$. Para el dise\~no
de redes tiene tres consecuencias, desarrolladas en el \Cref{ch:groups}: los
canales de una capa equivariante se indexan por irreducibles $\lambda$; el
n\'umero de par\'ametros por canal $d_\lambda^2$ queda fijado por el grupo, con
independencia de la resoluci\'on de entrada; y las convoluciones equivariantes
se convierten en multiplicaciones punto a punto en este dominio de Fourier
generalizado (espectral).

\section{Espacios homog\'eneos y cocientes}
\label{sec:found:homogeneous}

Hasta ahora las simetr\'ias han actuado sobre los datos. Ahora invertimos la
construcci\'on y usamos un grupo para \emph{construir} el propio dominio. Muchos
de los espacios en los que viven los datos geom\'etricos ---la esfera $S^2$, los
espacios proyectivos, las grassmannianas--- comparten una estructura oculta: son
cocientes $G/H$ de un grupo de Lie por un subgrupo cerrado. Reconocer un dominio
como tal cociente es el paso m\'as \'util al dise\~nar una arquitectura
equivariante para \'el.

Empezamos por la construcci\'on algebraica. Cuando $H \triangleleft \group$ es
normal, las clases laterales izquierdas $g H = \{gh : h \in H\}$ particionan
$\group$ y heredan una estructura de grupo.

\begin{definition}[Grupo cociente]
\label{def:found:quotient}
Sea $H \triangleleft \group$. El \emph{grupo cociente} $\group/H$ es el conjunto
de clases laterales izquierdas $\{gH : g \in \group\}$ con la operaci\'on
$(g_1 H)(g_2 H) = (g_1 g_2) H$.
\end{definition}

La operaci\'on est\'a bien definida precisamente porque $H$ es normal. El v\'inculo
entre cocientes y homomorfismos es el primer teorema de isomorf\'ia: para todo
homomorfismo $\phi : \group \to \mathcal{H}$,
\[
  \group / \ker(\phi) \;\cong\; \operatorname{Im}(\phi),
\]
mediante $g\,\ker(\phi) \mapsto \phi(g)$. Los cocientes modelan la reducci\'on
sistem\'atica de simetr\'ia que ocurre en arquitecturas jer\'arquicas: una red
sobre datos tridimensionales puede comenzar con la simetr\'ia completa $SE(3)$ y
reducirse a $\SO(3)$ en capas m\'as profundas mediante un homomorfismo cuyo
n\'ucleo (las traslaciones) es un subgrupo normal.

Para la geometr\'ia necesitamos dos nociones asociadas a una acci\'on: qu\'e
puntos puede alcanzar un punto dado y qu\'e elementos del grupo lo fijan.

\begin{definition}[\'Orbita y estabilizador]
\label{def:found:orbit}
Sea $\group$ actuando sobre un conjunto $X$ y sea $x \in X$. La \emph{\'orbita}
de $x$ es $\Orb(x) = \{g \cdot x : g \in \group\}$, el conjunto de puntos
alcanzables desde $x$. El \emph{estabilizador} de $x$ es
$\Stab(x) = \{g \in \group : g \cdot x = x\}$, el subgrupo que fija $x$.
\end{definition}

Ambas nociones se vinculan mediante el teorema \'orbita--estabilizador, que es el
coraz\'on geom\'etrico de esta secci\'on (\Cref{fig:found:orbit}).

\begin{theorem}[\'Orbita--estabilizador]
\label{thm:found:orbitstab}
Sea $\group$ actuando sobre $X$ y sea $x \in X$. Entonces existe una biyecci\'on
natural
\[
  \group / \Stab(x) \;\cong\; \Orb(x), \qquad g\,\Stab(x) \mapsto g \cdot x.
\]
\end{theorem}

\begin{proof}
La aplicaci\'on $\psi(g\,\Stab(x)) = g\cdot x$ est\'a bien definida: si
$g\,\Stab(x) = h\,\Stab(x)$ entonces $h^{-1}g \in \Stab(x)$, de donde
$g\cdot x = h\cdot((h^{-1}g)\cdot x) = h\cdot x$. Es inyectiva por el mismo
c\'alculo le\'ido al rev\'es, y suprayectiva porque todo punto de $\Orb(x)$ es
$g\cdot x$ para alg\'un $g$.
\end{proof}

\begin{figure}[h]
\centering
\begin{tikzpicture}[
    >=Stealth,
    pt/.style={circle, fill=#1, minimum size=7pt, inner sep=0pt}
]
% Grupo G con clases laterales
\draw[thick, fill=gdlgreen!12, rounded corners] (-5.2,-1.6) rectangle (-3.4,2.1);
\node[above, font=\small\bfseries, text=gdlgreen!60!black] at (-4.3,2.1) {Grupo $\group$};
\foreach \y/\lab in {1.3/$g_1\Stab(x)$, 0.3/$g_2\Stab(x)$, -0.7/$\Stab(x)$} {
    \draw[thick, fill=gdlgreen!22, rounded corners=1pt] (-5.05,\y) rectangle (-3.55,\y-0.55);
    \node[font=\tiny] at (-4.3,\y-0.27) {\lab};
}
% Flecha de biyecci\'on
\draw[->, very thick, gdlred!75!black] (-3.2,0.4) -- node[above, font=\small] {$\cong$} (-1.2,0.4);
% Espacio X con \'orbita
\draw[thick, fill=gdlgray!12, rounded corners] (-1,-1.6) rectangle (4.6,2.6);
\node[above, font=\small\bfseries, text=gdlgray!50!black] at (1.8,2.6) {Espacio $X$};
\draw[gdlblue!60, thick, fill=gdlblue!12] (1.7,0.5) ellipse (1.9cm and 0.95cm);
\node[above, font=\tiny, text=gdlblue!70!black] at (1.7,1.5) {$\Orb(x) = \group\cdot x$};
\node[pt=gdlred!75] (x) at (0.4,0.5) {};
\node[below left, font=\small] at (x) {$x$};
\node[pt=gdlorange!80] (gx1) at (2.2,1.1) {};
\node[above, font=\small] at (gx1) {$g_1\cdot x$};
\node[pt=gdlorange!80] (gx2) at (2.7,-0.1) {};
\node[below, font=\small] at (gx2) {$g_2\cdot x$};
\draw[->, thick, gdlpurple!75] (x) to[bend left=28] node[above, font=\tiny] {$g_1$} (gx1);
\draw[->, thick, gdlpurple!75] (x) to[bend right=18] node[below, font=\tiny] {$g_2$} (gx2);
\draw[->, thick, gdlgreen!55!black, loop above, looseness=6] (x) to node[above, font=\tiny] {$h\in\Stab(x)$} (x);
\end{tikzpicture}
\caption[Teorema \'orbita--estabilizador]{El teorema \'orbita--estabilizador.
Cada clase lateral $g\,\Stab(x)$ del estabilizador corresponde a exactamente un
punto $g\cdot x$ de la \'orbita, lo que da la biyecci\'on
$\group/\Stab(x) \cong \Orb(x)$.}
\label{fig:found:orbit}
\end{figure}

Cuando $\group$ es un grupo de Lie que act\'ua transitivamente y $H$ es un
subgrupo cerrado, $G/H$ es a su vez una variedad diferenciable, llamada
\emph{espacio homog\'eneo}, sobre la que $\group$ act\'ua transitivamente
mediante $g'\cdot(gH) = (g'g)H$. El teorema \'orbita--estabilizador afirma que
cualquier dominio que admita un grupo transitivo de simetr\'ias es un espacio
homog\'eneo: se identifica el grupo de simetr\'ia $\group$, se fija un punto de
referencia, y el estabilizador $H$ de ese punto determina el dominio como $G/H$.

\begin{example}[La esfera como $\SO(3)/\SO(2)$]
\label{ex:found:sphere}
Sea $\SO(3)$ actuando sobre $S^2$ por rotaci\'on. La acci\'on es transitiva:
cualquier vector unitario puede rotarse a cualquier otro. Fijemos el polo norte
$p = (0,0,1)^\top$; una rotaci\'on fija $p$ exactamente cuando su eje es el eje
$z$, y esas rotaciones forman un subgrupo isomorfo a $\SO(2)$. Por tanto
$\Stab(p) \cong \SO(2)$, y el teorema \'orbita--estabilizador da
\[
  S^2 \;\cong\; \SO(3)/\SO(2).
\]
Una se\~nal $f : S^2 \to \R^C$ se eleva a una funci\'on sobre $\SO(3)$ que es
invariante a derecha bajo $\SO(2)$, a saber $\tilde f(g) = f(g\cdot p)$. En
consecuencia, la convoluci\'on equivariante natural sobre la esfera es una
correlaci\'on sobre el grupo completo $\SO(3)$ en lugar de sobre una proyecci\'on
plana, que es la construcci\'on que subyace a las redes esf\'ericas discutidas en
los \Cref{ch:geodesics,ch:applications}.
\end{example}

Este ejemplo ilustra una receta general para arquitecturas geom\'etricas: (1)
identificar el dominio de los datos $\Omega$; (2) determinar el grupo de
simetr\'ia relevante $\group$; (3) calcular el estabilizador $H$ de un punto de
referencia; (4) reconocer $\Omega \cong G/H$ por el teorema
\'orbita--estabilizador; y (5) construir la arquitectura equivariante usando el
an\'alisis arm\'onico de $\group$ (\Cref{thm:found:peterweyl}). Los espacios
homog\'eneos disfrutan adem\'as de propiedades geom\'etricas convenientes: si
$\group$ es conexo entonces $G/H$ es conexo, y $\group$ act\'ua por isometr\'ias
de cualquier m\'etrica riemanniana $\group$-invariante sobre $G/H$. De forma
crucial, sobre un espacio homog\'eneo un operador equivariante queda
completamente determinado por su comportamiento en el punto base $o = eH$: como
$\group$ act\'ua transitivamente, el valor de $\Phi f$ en cualquier punto
$x = g\cdot o$ queda fijado por la equivarianza, y un breve c\'alculo muestra que
los \'unicos operadores lineales equivariantes son convoluciones generalizadas.
Esta es la raz\'on matem\'atica por la que la convoluci\'on es la operaci\'on
can\'onica sobre dominios sim\'etricos, un hilo que seguimos a lo largo de los
\Cref{ch:grids,ch:groups,ch:graphs}.

\section{Variedades y espacios tangentes}
\label{sec:found:manifolds}

El tercer hilo de nuestro lenguaje es la geometr\'ia de los dominios curvos. Una
superficie, un espacio de configuraciones o el conjunto de rotaciones no es un
espacio vectorial, y sin embargo cada uno se parece al espacio eucl\'ideo plano
cuando se mira de cerca. La noci\'on de variedad lo precisa, y el espacio
tangente proporciona la linearizaci\'on sobre la que procede el c\'alculo ---y, en
\'ultima instancia, el aprendizaje---.

Asumimos los prerrequisitos topol\'ogicos est\'andar: un \emph{espacio
topol\'ogico} es un conjunto con una colecci\'on distinguida de abiertos cerrada
bajo uniones arbitrarias e intersecciones finitas; es \emph{de Hausdorff} si
puntos distintos tienen entornos abiertos disjuntos; y un \emph{homeomorfismo} es
una biyecci\'on continua con inversa continua, la noci\'on de que dos espacios
son ``topol\'ogicamente iguales''.

\begin{definition}[Variedad topol\'ogica]
\label{def:found:topmanifold}
Una \emph{variedad topol\'ogica} de dimensi\'on $m$ es un espacio topol\'ogico de
Hausdorff y con base numerable $\mathcal{M}$ que es \emph{localmente
euclidiano}: todo punto tiene un entorno abierto $U$ homeomorfo, mediante una
aplicaci\'on $\phi : U \to V \subseteq \R^m$ (una \emph{carta}), a un abierto de
$\R^m$.
\end{definition}

Intuitivamente, una variedad es un espacio que localmente se parece a $\R^m$
aunque su forma global sea muy distinta: la superficie de la Tierra es curva
pero parece plana a un peat\'on; una hoja de papel arrugada es localmente
bidimensional; el espacio de configuraciones de un brazo rob\'otico es no
eucl\'ideo pero localmente se coordina con los \'angulos de las articulaciones.
Una colecci\'on de cartas que cubre $\mathcal{M}$ es un \emph{atlas}
$\mathcal{A} = \{(U_\alpha, \phi_\alpha)\}$, y las \emph{funciones de
transici\'on} $\phi_\beta \circ \phi_\alpha^{-1}$ describen c\'omo se relacionan
las cartas solapadas. La variedad es \emph{diferenciable} (suave) si todas las
funciones de transici\'on son $C^\infty$; esta es exactamente la estructura
adicional necesaria para hacer c\'alculo sobre $\mathcal{M}$.

El objeto central es el espacio tangente, que captura las direcciones
infinitesimales en un punto. Dos construcciones complementarias, ilustradas en
la \Cref{fig:found:tangent}, lo precisan.

\begin{definition}[Espacio tangente v\'ia curvas]
\label{def:found:tangent-curves}
Sea $\mathcal{M}$ suave y $p \in \mathcal{M}$. Dos curvas
$\gamma_1, \gamma_2 : (-\epsilon,\epsilon) \to \mathcal{M}$ con
$\gamma_1(0) = \gamma_2(0) = p$ son \emph{tangentes en $p$} si, para alguna (y
por tanto toda) carta $(U,\phi)$ alrededor de $p$,
$(\phi\circ\gamma_1)'(0) = (\phi\circ\gamma_2)'(0)$. El \emph{espacio tangente}
$T_p\mathcal{M}$ es el conjunto de clases de equivalencia $[\gamma]$ de curvas
tangentes en $p$.
\end{definition}

\begin{definition}[Espacio tangente v\'ia derivaciones]
\label{def:found:tangent-deriv}
Una \emph{derivaci\'on} en $p$ es una aplicaci\'on lineal
$v : C^\infty(\mathcal{M}) \to \R$ que obedece la regla de Leibniz
$v(fg) = v(f)\,g(p) + f(p)\,v(g)$. El espacio tangente $T_p\mathcal{M}$ es el
espacio vectorial de todas las derivaciones en $p$.
\end{definition}

La primera construcci\'on lee un vector tangente como la velocidad inicial de una
curva ---una direcci\'on de movimiento---; la segunda lo lee como un operador de
derivada direccional ---una direcci\'on de cambio de funciones---. Ambas son
can\'onicamente isomorfas, enviando el isomorfismo $[\gamma]$ a la derivaci\'on
$f \mapsto \frac{d}{dt}\big|_0 (f\circ\gamma)$. En una carta con coordenadas
$x^1,\dots,x^m$ las derivadas parciales
$\frac{\partial}{\partial x^j}\big|_p$ forman una base, de modo que
$T_p\mathcal{M}$ es un espacio vectorial de dimensi\'on $m = \dim\mathcal{M}$.
Usamos curvas para la intuici\'on geom\'etrica y derivaciones para la
manipulaci\'on algebraica.

\begin{example}[Espacio tangente a la esfera]
\label{ex:found:tangent-sphere}
En el polo norte $p = (0,0,1)$ de $S^2$, la curva
$\gamma(t) = (\sin t, 0, \cos t)$ tiene $\gamma'(0) = (1,0,0)$. El espacio
tangente $T_p S^2$ es el plano horizontal $\{(x,y,0)\}$, el plano tangente a la
esfera en el polo. Su bidimensionalidad refleja $\dim S^2 = 2$.
\end{example}

\begin{figure}[h]
\centering
\begin{tikzpicture}[>=Stealth, scale=1.0]
% Superficie
\shade[ball color=gdlblue!18, opacity=0.9]
    (-2,-1.3) .. controls (-1,1) and (1,2) .. (2.2,0.6)
    .. controls (1.6,-1) and (-1,-1.9) .. (-2,-1.3);
\draw[thick, gdlblue!55]
    (-2,-1.3) .. controls (-1,1) and (1,2) .. (2.2,0.6)
    .. controls (1.6,-1) and (-1,-1.9) .. (-2,-1.3);
\node[text=gdlblue!60!black, font=\small] at (-1.7,-1.5) {$\mathcal{M}$};
% Punto p
\coordinate (p) at (0,0.35);
% Plano tangente (paralelogramo)
\begin{scope}[on background layer]
\fill[gdlorange!18, opacity=0.85]
    ($(p)+(-1.7,-0.55)$) -- ($(p)+(1.7,0.35)$) --
    ($(p)+(1.4,1.15)$) -- ($(p)+(-2.0,0.25)$) -- cycle;
\draw[gdlorange!60]
    ($(p)+(-1.7,-0.55)$) -- ($(p)+(1.7,0.35)$) --
    ($(p)+(1.4,1.15)$) -- ($(p)+(-2.0,0.25)$) -- cycle;
\end{scope}
\node[text=gdlorange!60!black, font=\small] at (1.9,1.25) {$T_p\mathcal{M}$};
% Curvas por p
\draw[thick, gdlpurple!75] (-1.4,-0.4) .. controls (-0.5,0.05) .. (p)
    .. controls (0.5,0.65) .. (1.4,1.0);
\node[gdlpurple!75, font=\tiny] at (1.55,1.05) {$\gamma_1$};
\draw[thick, gdlgreen!60!black] (-0.5,-0.9) .. controls (-0.2,-0.25) .. (p)
    .. controls (0.2,0.9) .. (0.45,1.4);
\node[gdlgreen!60!black, font=\tiny] at (0.6,1.45) {$\gamma_2$};
% Vectores tangentes
\draw[->, very thick, gdlpurple!75] (p) -- ++(0.95,0.45)
    node[right, font=\tiny] {$\gamma_1'(0)$};
\draw[->, very thick, gdlgreen!60!black] (p) -- ++(0.25,0.85)
    node[above, font=\tiny] {$\gamma_2'(0)$};
\fill[gdlred!80] (p) circle (2.2pt);
\node[below right=0pt and 1pt of p, font=\small] {$p$};
\end{tikzpicture}
\caption[Espacio tangente]{El espacio tangente $T_p\mathcal{M}$ en un punto $p$.
Los vectores tangentes son las velocidades iniciales $\gamma'(0)$ de curvas que
pasan por $p$; de forma equivalente, operadores de derivada direccional sobre
funciones suaves. Forman un espacio vectorial que aproxima $\mathcal{M}$
linealmente en $p$.}
\label{fig:found:tangent}
\end{figure}

Una aplicaci\'on suave $F : \mathcal{M} \to \mathcal{N}$ act\'ua sobre los
vectores tangentes mediante su \emph{diferencial} (o pushforward)
$dF_p : T_p\mathcal{M} \to T_{F(p)}\mathcal{N}$, definida sobre derivaciones por
$dF_p(v)(g) = v(g\circ F)$; en coordenadas locales $dF_p$ es la matriz
jacobiana. El diferencial clasifica las aplicaciones por el comportamiento de su
rango: $F$ es una \emph{inmersi\'on} si $dF_p$ es inyectiva en todo punto y una
\emph{submersi\'on} si es suprayectiva en todo punto. Un \emph{embebimiento} es
una inmersi\'on inyectiva que es un homeomorfismo sobre su imagen;
intuitivamente, una variedad embebida se sit\'ua dentro del espacio ambiente sin
autointersecciones ni picos. Estas nociones reaparecen en el
\Cref{ch:geodesics}, donde una m\'etrica riemanniana sobre los espacios
tangentes convierte una variedad en un espacio con distancias y curvatura, y en
el \Cref{ch:gauges}, donde la ausencia de una base can\'onica para
$T_p\mathcal{M}$ conduce a la simetr\'ia gauge.

Conectando con la \Cref{sec:found:groups}: un grupo de Lie es una variedad cuyo
espacio tangente en la identidad es su \'algebra de Lie, y el mapa exponencial es
la curva can\'onica a trav\'es de la identidad en cada direcci\'on. Los tres
hilos de este cap\'itulo ---grupos, representaciones y variedades--- confluyen
as\'i en el grupo de Lie, y juntos proporcionan el vocabulario para los sesgos
geom\'etricos del \Cref{ch:priors} y las arquitecturas que les siguen.
